% ============================================================
%  PRESENTACIÓN: Artículo Científico y Ética Académica
%  Basado en el template Obsidian & Ember
% ============================================================
\documentclass[aspectratio=169, 12pt]{beamer}

\usepackage{fontspec}
\usepackage{tikz}
\usepackage{xcolor}
\usepackage{graphicx}
\usepackage{microtype}
\usepackage{tcolorbox}
\usepackage{calc}
\usepackage{ragged2e}

\usetikzlibrary{positioning, shapes.geometric, fadings}

% ── Fonts ────────────────────────────────────────────────────
\setmainfont{Times New Roman}
\setsansfont{Arial}
\setmonofont{Consolas}[Scale=MatchLowercase]

% ── Colour palette ───────────────────────────────────────────
\definecolor{Obsidian}{HTML}{0D0D0D}
\definecolor{Charcoal}{HTML}{1A1A1A}
\definecolor{Slate}{HTML}{2C2C2C}
\definecolor{Ember}{HTML}{E8643A}
\definecolor{Gold}{HTML}{C49A3C}
\definecolor{Smoke}{HTML}{8A8A8A}
\definecolor{Cloud}{HTML}{D6D6D6}
\definecolor{OffWhite}{HTML}{F5F5F5}
\definecolor{SteelBlue}{HTML}{5B8FA8}

% ── Base Beamer settings ─────────────────────────────────────
\usetheme{default}
\usecolortheme{default}
\setbeamertemplate{navigation symbols}{}
\setbeamertemplate{footline}{}

\setbeamercolor{background canvas}{bg=Obsidian}
\setbeamercolor{normal text}{fg=Cloud}
\setbeamercolor{frametitle}{fg=OffWhite}
\setbeamertemplate{frametitle}{\vspace{8pt}{\color{OffWhite}\rmfamily\bfseries\Large\insertframetitle}\ifx\insertframesubtitle\empty\else\\[2pt]{\color{Smoke}\sffamily\small\insertframesubtitle}\fi\vspace{4pt}\emberrule}
\setbeamerfont{normal text}{family=\sffamily}

% Remove default beamer chrome
\setbeamertemplate{headline}{}
\setbeamertemplate{footline}{}
\setbeamertemplate{sidebar left}{}
\setbeamertemplate{sidebar right}{}

% ── Reusable helpers ─────────────────────────────────────────

% Ember accent rule (full-width top bar variant)
\newcommand{\emberrule}{%
  
\begin{tikzpicture}
    \fill[Ember] (0,0) rectangle (\textwidth, 1.5pt);
  \end{tikzpicture}\vspace{4pt}%
}

% Chip / label badge
\newcommand{\chip}[2][Ember]{%
  \tcbox[on line, colback=#1, colframe=#1, arc=1pt,
         boxsep=1pt, left=4pt, right=4pt, top=1.5pt, bottom=1.5pt,
         fontupper=\sffamily\tiny\bfseries\color{OffWhite}\MakeUppercase]{#2}%
}

% Stat card (number + label)
\newcommand{\statcard}[3][Ember]{%
  \begin{tcolorbox}[
    colback=Slate, colframe=#1, arc=0pt,
    leftrule=4pt, rightrule=0pt, toprule=0pt, bottomrule=0pt,
    left=6pt, right=6pt, top=6pt, bottom=6pt,
  ]
    {\color{OffWhite}\rmfamily\bfseries\fontsize{26}{28}\selectfont #2}\\[3pt]
    {\color{Smoke}\sffamily\fontsize{6.5}{6.5}\selectfont\MakeUppercase{#3}}
  \end{tcolorbox}%
}

% Image placeholder (replace \imageph with \includegraphics in real use)
\newcommand{\imageph}[2][5cm]{%
  \begin{tikzpicture}
    \fill[Charcoal] (0,0) rectangle (\linewidth,#1);
    \node[text=Slate!60!Charcoal, font=\sffamily\tiny\itshape]
      at (0.5\linewidth, 0.5*#1) {#2};
  \end{tikzpicture}%
}

% Icon illustrations (TikZ) to avoid external images
\newcommand{\bookicon}[1][2cm]{%
  \begin{tikzpicture}[scale=1, baseline={(current bounding box.center)}]
    \fill[OffWhite] (0,0) rectangle (#1,0.8#1);
    \draw[black,line width=0.8pt] (0,0) rectangle (#1,0.8#1);
    \fill[Slate] (0.05#1,0.1#1) rectangle (0.25#1,0.7#1);
    \draw (0.03#1,0.1#1) -- (0.03#1,0.7#1);
    \draw (0.27#1,0.1#1) -- (0.27#1,0.7#1);
    \node[anchor=north west, font=\sffamily\tiny] at (0.3#1,0.7#1) {ARTÍCULO};
  \end{tikzpicture}%
}

\newcommand{\mapicon}[1][2cm]{%
  \begin{tikzpicture}[scale=1, baseline={(current bounding box.center)}]
    \fill[OffWhite] (0,0) rectangle (#1,0.8#1);
    \draw[black,line width=0.8pt] (0,0) rectangle (#1,0.8#1);
    \draw[Slate,line width=0.9pt] (0.08#1,0.6#1) -- (0.3#1,0.72#1) -- (0.55#1,0.45#1) -- (0.78#1,0.6#1);
    \node[anchor=south, font=\sffamily\tiny, text=Ember] at (0.35#1,0.2#1) {C.R.};
  \end{tikzpicture}%
}

\newcommand{\balanceicon}[1][2cm]{%
  \begin{tikzpicture}[scale=1, baseline={(current bounding box.center)}]
    \draw[black,line width=0.8pt] (0.5,0) -- (0.5,0.8#1);
    \draw[black,line width=0.8pt] (0.2,0.8#1) -- (0.8,0.8#1);
    \draw[black,line width=0.6pt] (0.25,0.8#1) -- (0.5,0.7#1) -- (0.75,0.8#1);
    \draw[black,line width=0.6pt] (0.28,0.52#1) -- (0.31,0.2#1); 
    \draw[black,line width=0.6pt] (0.72,0.52#1) -- (0.69,0.2#1);
    \fill[Ember] (0.22,0.28#1) circle (0.03#1);
    \fill[Ember] (0.68,0.28#1) circle (0.03#1);
  \end{tikzpicture}%
}

% Full-bleed background helper (used inside tikzpicture overlays)
% Usage: \fullbg{color} inside frame with [plain]
\newcommand{\fullbg}[1]{%
  \begin{tikzpicture}[remember picture, overlay]
    \fill[#1] (current page.south west) rectangle (current page.north east);
  \end{tikzpicture}%
}

% ── Footer bar ───────────────────────────────────────────────
\setbeamertemplate{footline}{%
  \begin{tikzpicture}[remember picture, overlay]
    \fill[Charcoal]
      (current page.south west) rectangle
      ([yshift=14pt]current page.south east);
    \node[anchor=west, text=Smoke, font=\sffamily\fontsize{5}{5}\selectfont,
          inner sep=0] at ([xshift=10pt, yshift=7pt]current page.south west)
      {OBSIDIAN \& EMBER TEMPLATE};
    \node[anchor=east, text=Smoke, font=\sffamily\fontsize{5}{5}\selectfont,
          inner sep=0] at ([xshift=-10pt, yshift=7pt]current page.south east)
      {\insertframenumber\,/\,\inserttotalframenumber};
  \end{tikzpicture}%
}

% ============================================================
\begin{document}
% ============================================================

% ── SLIDE 1: Title ───────────────────────────────────────────
% Sugerencia: Reemplazar [ HERO IMAGE ] con \includegraphics{imagen_investigacion.jpg} (ej. libro abierto o microscopio)
{
\setbeamertemplate{footline}{}   % no footer on title
\begin{frame}[plain]
  \begin{tikzpicture}[remember picture, overlay]

    % Left image panel
    \fill[Charcoal]
      (current page.south west) rectangle
      ([xshift=0.52\paperwidth]current page.north west);
    % Ember vertical rule
    \fill[Ember]
      ([xshift=0.52\paperwidth]current page.south west)
      rectangle
      ([xshift=0.525\paperwidth]current page.north west);

    % Image placeholder text
    \node[text=Slate!60!Charcoal, font=\sffamily\tiny\itshape]
      at ([xshift=0.1\paperwidth]current page.center)
      {\includegraphics[width=0.4\paperwidth]{imagenes/portada.png}};

    % Eyebrow label
    \node[anchor=north west,
          text=Gold,
          font=\sffamily\fontsize{7}{7}\selectfont\bfseries,
          inner sep=0,
          xshift=0.55\paperwidth, yshift=-0.10\paperheight]
      at (current page.north west)
      {\MakeUppercase{Maestría en Ciencias de la Computación}};

    % Big title — word 1
    \node[anchor=north west,
          text=OffWhite,
          font=\rmfamily\fontsize{32}{32}\selectfont\bfseries,
          inner sep=0,
          xshift=0.545\paperwidth, yshift=-0.20\paperheight]
      at (current page.north west) {ARTÍCULO};

    % Big title — word 2 (accent)
    \node[anchor=north west,
          text=Ember,
          font=\rmfamily\fontsize{32}{32}\selectfont\bfseries,
          inner sep=0, text width=0.42\paperwidth,
          xshift=0.545\paperwidth, yshift=-0.40\paperheight]
      at (current page.north west) {CIENTÍFICO};

    % Subtitle
    \node[anchor=north west,
          text=Cloud,
          font=\sffamily\fontsize{9}{13}\selectfont,
          text width=0.39\paperwidth,
          inner sep=0,
          xshift=0.55\paperwidth, yshift=-0.60\paperheight]
      at (current page.north west)
      {y Ética Académica en\\ Costa Rica.};

    % Divider line
    \draw[Slate, line width=0.4pt]
      ([xshift=0.55\paperwidth, yshift=0.18\paperheight]current page.south west)
      -- ([xshift=0.97\paperwidth, yshift=0.18\paperheight]current page.south west);

    % Presenter info
    \node[anchor=west,
          text=Smoke,
          font=\sffamily\fontsize{6.5}{6.5}\selectfont,
          inner sep=0,
          xshift=0.55\paperwidth, yshift=0.14\paperheight]
      at (current page.south west)
      {EMMANUEL BARRANTES  ·  ABRIL 2026};

  \end{tikzpicture}
\end{frame}
}

% ── SLIDE 2: Mapa de la Presentación ─────────────────────────
\begin{frame}{Mapa de la Presentación}
  \framesubtitle{15 minutos divididos en tres bloques inseparables}
  \vspace{-4pt}
  \begin{columns}[T, totalwidth=\textwidth]

    %% Bloque 1
    \begin{column}{0.315\textwidth}
      \begin{tcolorbox}[
        colback=Slate, colframe=Slate, arc=1pt,
        left=0pt, right=0pt, top=0pt, bottom=5pt,
        toprule=2pt, colframe=Ember,
      ]
        \centerline{\bookicon[2.1cm]}\\[4pt]
        \hspace{5pt}\chip{1:00 - 6:00}\\[4pt]
        \hspace{5pt}{\color{OffWhite}\rmfamily\bfseries\normalsize El Artículo\\ Científico}\\[4pt]
        \hspace{5pt}{\color{Cloud}\sffamily\fontsize{7.5}{11}\selectfont
          ¿Qué es y qué no es?\\ Tres preguntas clave.\\ Abstract y tipos.}\\[4pt]
        \hspace{5pt}{\color{Ember}\sffamily\fontsize{7}{7}\selectfont\bfseries
          Rigor \textrightarrow}
      \end{tcolorbox}
    \end{column}

    \hfill

    %% Bloque 2
    \begin{column}{0.315\textwidth}
      \begin{tcolorbox}[
        colback=Slate, colframe=Slate, arc=1pt,
        left=0pt, right=0pt, top=0pt, bottom=5pt,
        toprule=2pt, colframe=Gold,
      ]
        \centerline{\mapicon[2.1cm]}\\[4pt]
        \hspace{5pt}\chip[Gold]{6:00 - 9:00}\\[4pt]
        \hspace{5pt}{\color{OffWhite}\rmfamily\bfseries\normalsize Dónde Publicar\\ en Costa Rica}\\[4pt]
        \hspace{5pt}{\color{Cloud}\sffamily\fontsize{7.5}{11}\selectfont
          Revistas y congresos.\\ JoCICI, CLEI, Revista Ingeniería.}\\[4pt]
        \hspace{5pt}{\color{Gold}\sffamily\fontsize{7}{7}\selectfont\bfseries
          Estrategia \textrightarrow}
      \end{tcolorbox}
    \end{column}

    \hfill

    %% Bloque 3
    \begin{column}{0.315\textwidth}
      \begin{tcolorbox}[
        colback=Slate, colframe=Slate, arc=1pt,
        left=0pt, right=0pt, top=0pt, bottom=5pt,
        toprule=2pt, colframe=SteelBlue,
      ]
        \centerline{\balanceicon[2.1cm]}\\[4pt]
        \hspace{5pt}\chip[SteelBlue]{9:00 - 13:00}\\[4pt]
        \hspace{5pt}{\color{OffWhite}\rmfamily\bfseries\normalsize Ética\\ Académica}\\[4pt]
        \hspace{5pt}{\color{Cloud}\sffamily\fontsize{7.5}{11}\selectfont
          Plagio, Ley 6683.\\ Sanciones en universidades.}\\[4pt]
        \hspace{5pt}{\color{SteelBlue}\sffamily\fontsize{7}{7}\selectfont\bfseries
          Honestidad \textrightarrow}
      \end{tcolorbox}
    \end{column}

  \end{columns}
\end{frame}

% ── SLIDE 3: ¿Qué NO es un Artículo? ─────────────────────────
% Sugerencia: Incluir ícono de artículo científico o pluma tachada
\begin{frame}{¿Qué NO es un Artículo?}
  \framesubtitle{Definición en negativo: una sola idea y evidencia}
  \vspace{-2pt}

  \begin{itemize}
    \item No es el manual de la herramienta que desarrollaste.
    \item No es una narración de todo lo que hiciste durante dos años.
    \item No es propaganda de tu proyecto.
  \end{itemize}

  \vspace{0.5cm}
  {\color{OffWhite}\rmfamily\bfseries\normalsize Un artículo es la comunicación objetiva de una sola idea y la evidencia que la respalda.}
  \vspace{0.3cm}
  {\color{Cloud}\sffamily\fontsize{8}{10}\selectfont Esa restricción — entre ocho y doce páginas — es lo que lo hace difícil.}
\end{frame}

% ── SLIDE 4: Tres Preguntas Obligatorias ─────────────────────
\begin{frame}{Tres Preguntas Obligatorias}
  \framesubtitle{De Snyder (1991): antes de escribir una línea}
  \vspace{-2pt}

  \begin{enumerate}
    \item ¿Por qué escribo este artículo? (Si es 'para documentar lo que hice', es un informe técnico.)
    \item ¿Qué estoy comunicando? (Debe resumirse en una oración.)
    \item ¿A quién va dirigido? (Define el vehículo, nivel técnico y términos.)
  \end{enumerate}

  \vspace{0.5cm}
  {\color{Ember}\sffamily\bfseries Si no puedes responder con precisión alguna, el artículo aún no está listo.}
\end{frame}

% ── SLIDE 5: El Abstract: Spoiler, no Tráiler ─────────────────
\begin{frame}{El Abstract: Spoiler, no Tráiler}
  \framesubtitle{Debe declarar el resultado desde el inicio}
  \vspace{-2pt}

  \begin{columns}[T]
    \begin{column}{0.48\textwidth}
      {\color{Smoke}\sffamily\small\bfseries Abstract Débil:}\\[2pt]
      {\color{Cloud}\sffamily Este trabajo estudia las redes neuronales aplicadas a X.}
    \end{column}
    \hfill
    \begin{column}{0.48\textwidth}
      {\color{Gold}\sffamily\small\bfseries Abstract con Contribución:}\\[2pt]
      {\color{Cloud}\sffamily Este trabajo demuestra que el método Y mejora la precisión en un 20\% sobre el estado del arte.}
    \end{column}
  \end{columns}

  \vspace{0.5cm}
  {\color{OffWhite}\rmfamily\bfseries Si la contribución no cabe en un párrafo, hay un problema de fondo.}
\end{frame}

% ── SLIDE 6: Tipos Principales de Artículos ───────────────────
\begin{frame}{Tipos Principales de Artículos}
  \framesubtitle{En Ciencias de la Computación}
  \vspace{-2pt}

  \begin{itemize}
    \item {\bfseries Teórico:} Definiciones + pruebas lógicas formales.
    \item {\bfseries Reporte de Experiencia:} Generalización de observaciones, no narración cronológica.
    \item {\bfseries Sobre Métodos:} Comparación objetiva con métricas claras + reconocimiento de limitaciones.
  \end{itemize}
\end{frame}

% ── SLIDE 7: Dónde Publicar en Costa Rica ────────────────────
% Sugerencia: Incluir mapa de Costa Rica o ícono de publicación
\begin{frame}{Dónde Publicar en Costa Rica}
  \framesubtitle{Diferencia con Brasil: orientación directa a índices internacionales}
  \vspace{-2pt}

  {\color{Cloud}\sffamily Costa Rica no tiene un clasificador nacional como Qualis-Capes. Los programas de posgrado se orientan por Scopus, WoS y SciELO.}
\end{frame}

% ── SLIDE 8: Revistas Nacionales ─────────────────────────────
% Sugerencia: Incluir logos de Revista Ingeniería, Repertorio Científico, Investiga.TEC
\begin{frame}{Revistas Nacionales}
  \framesubtitle{Principales opciones para Computación}
  \vspace{-2pt}

  \begin{itemize}
    \item {\bfseries Revista Ingeniería (UCR):} Más relevante. Desde 1991, INII-UCR. Revisión doble ciego. Indexada en SciELO, Latindex, DOAJ.
    \item {\bfseries Repertorio Científico (UNED):} Cobertura en Informática. Acceso abierto, sin costo para autores.
    \item {\bfseries Investiga.TEC:} Carácter divulgativo, no revisión científica estricta.
  \end{itemize}
\end{frame}

% ── SLIDE 9: Congresos de Referencia ─────────────────────────
\begin{frame}{Congresos de Referencia}
  \framesubtitle{JoCICI y CLEI: eventos clave}
  \vspace{-2pt}

  \begin{itemize}
    \item {\bfseries JoCICI:} Jornadas Costarricenses (2013-). Formato IEEE, 6 páginas. Proceedings a IEEE Xplore. Edición 2026 en UCR.
    \item {\bfseries CLEI:} Conferencia Latinoamericana (1974-). Itinerante. Mejores trabajos a CLEIej (Latindex, DOAJ). Costa Rica sede 2021.
  \end{itemize}
\end{frame}

% ── SLIDE 10: Tabla Resumen de Vehículos ─────────────────────
\begin{frame}{Tabla Resumen de Vehículos}
  \framesubtitle{Comparación rápida}
  \vspace{-2pt}

  \begin{tabular}{|l|l|l|l|}
    \hline
    Vehículo & Tipo & Alcance & Indexación \\
    \hline
    Revista Ingeniería (UCR) & Revista & Nacional-regional & SciELO, Latindex, DOAJ \\
    Repertorio Científico (UNED) & Revista & Nacional & Latindex \\
    JoCICI & Congreso & Nacional & IEEE Xplore \\
    CLEI + CLEIej & Congreso + Revista & Latinoamericano & Latindex, DOAJ \\
    LACLO & Congreso & Latinoamericano & - \\
    \hline
  \end{tabular}
\end{frame}

% ── SLIDE 11: UCRIndex y Accensum ────────────────────────────
\begin{frame}{UCRIndex y Accensum}
  \framesubtitle{Herramientas para navegar el ecosistema}
  \vspace{-2pt}

  \begin{itemize}
    \item {\bfseries UCRIndex:} Evalúa revistas UCR con 5 indicadores. Distribuye en cuartiles (A = mayor calidad).
    \item {\bfseries Accensum:} Lista revistas costarricenses con estado en SciELO CR, WoS y Scopus (actualizado marzo 2026).
  \end{itemize}
\end{frame}

% ── SLIDE 12: Plagio: Dos Formas ─────────────────────────────
% Sugerencia: Íconos de texto copiado (rojo) vs. texto original (verde)
\begin{frame}{Plagio: Dos Formas}
  \framesubtitle{Un mismo problema}
  \vspace{-2pt}

  \begin{itemize}
    \item {\bfseries Copia Literal:} Texto ajeno sin comillas ni fuente.
    \item {\bfseries Copia de Ideas:} Mismas ideas en secuencia lógica, aunque parafraseadas.
  \end{itemize}

  \vspace{0.5cm}
  {\color{Ember}\sffamily\bfseries Regla básica inamovible: comillas + fuente, siempre, sin excepción.}
\end{frame}

% ── SLIDE 13: Marco Legal en Costa Rica ──────────────────────
\begin{frame}{Marco Legal en Costa Rica}
  \framesubtitle{Ley No. 6683: plagio genera responsabilidad civil, no penal}
  \vspace{-2pt}

  \begin{itemize}
    \item Protección automática desde la creación (Reglamento 24611-J, Art. 2).
    \item Duración: vida del autor + 70 años (Art. 58).
    \item Excepción académica (Art. 73): usos honrados con fines educativos, parcialmente y con fuente.
  \end{itemize}
\end{frame}

% ── SLIDE 14: Sanciones en Universidades ─────────────────────
\begin{frame}{Sanciones en Universidades}
  \framesubtitle{Públicas costarricenses (Scielo CR, 2022)}
  \vspace{-2pt}

  \begin{itemize}
    \item {\bfseries UCR:} Falta grave (obra sin protección) / muy grave (obra protegida). Personal académico desde 2014.
    \item {\bfseries TEC:} 'Conducta fraudulenta'; sanciones progresivas.
    \item {\bfseries UNED:} Para posgrado, cualquier plagio es falta muy grave.
  \end{itemize}
\end{frame}

% ── SLIDE 15: Autoplagio y Envíos Simultáneos ────────────────
\begin{frame}{Autoplagio y Envíos Simultáneos}
  \framesubtitle{Dos reglas finales}
  \vspace{-2pt}

  \begin{itemize}
    \item {\bfseries Autoplagio:} Un artículo, un vehículo a la vez. Válido reutilizar texto propio de monografía origen.
    \item {\bfseries Envíos Simultáneos:} No múltiples versiones del mismo artículo.
  \end{itemize}
\end{frame}

% ── SLIDE 16: Lista de Verificación ───────────────────────────
\begin{frame}{Lista de Verificación}
  \framesubtitle{Antes de enviar}
  \vspace{-2pt}

  \begin{itemize}
    \item ¿Puedo resumir la contribución en una oración?
    \item ¿El abstract declara el resultado, no solo el tema?
    \item ¿Comparo con el estado del arte actual?
    \item ¿Menciono trabajos relacionados con comparación objetiva?
    \item ¿Todo texto ajeno está entre comillas y fuente?
    \item ¿El artículo va a un único vehículo?
    \item ¿El vehículo es adecuado al nivel de la contribución?
  \end{itemize}
\end{frame}

% ── SLIDE 17: Cierre ──────────────────────────────────────────
{
\setbeamertemplate{footline}{}
\begin{frame}[plain]
  
\begin{tikzpicture}[remember picture, overlay]

    % Ember left panel
    \fill[Ember]
      (current page.south west) rectangle
      ([xshift=0.38\paperwidth]current page.north west);

    % Decorative circles on left panel
    \foreach \cx/\cy in {0.08/0.85, 0.22/0.85, 0.08/0.60, 0.22/0.60,
                          0.08/0.35, 0.22/0.35, 0.08/0.10, 0.22/0.10} {
      \fill[black, opacity=0.15]
        ([xshift=\cx*\paperwidth, yshift=\cy*\paperheight]current page.south west)
        circle (0.048*\paperwidth);
    }

    % "THANK YOU" watermark on left panel
    \node[text=black, opacity=0.18,
          font=\rmfamily\bfseries\fontsize{42}{42}\selectfont,
          rotate=0,
          anchor=center]
      at ([xshift=0.19\paperwidth]current page.center)
      {\shortstack{GRACIAS}};

    % Right: headline
    \node[anchor=north west,
          text=OffWhite,
          font=\rmfamily\bfseries\fontsize{28}{32}\selectfont,
          inner sep=0,
          xshift=0.43\paperwidth, yshift=-0.15\paperheight]
      at (current page.north west) {Publicar con Rigor};

    % Subtitle
    \node[anchor=north west,
          text=Cloud,
          font=\sffamily\fontsize{9}{14}\selectfont,
          text width=0.50\paperwidth,
          inner sep=0,
          xshift=0.43\paperwidth, yshift=-0.34\paperheight]
      at (current page.north west)
      {y Honestidad no son\\ objetivos separados.\\ Son la misma cosa.};

    % CTA button
    \fill[Ember]
      ([xshift=0.43\paperwidth, yshift=0.39\paperheight]current page.south west)
      rectangle
      ([xshift=0.64\paperwidth, yshift=0.47\paperheight]current page.south west);
    \node[text=OffWhite,
          font=\sffamily\bfseries\fontsize{8}{8}\selectfont,
          anchor=center]
      at ([xshift=0.535\paperwidth, yshift=0.43\paperheight]current page.south west)
      {PREGUNTAS?};

    % Divider
    \draw[Slate, line width=0.4pt]
      ([xshift=0.43\paperwidth, yshift=0.34\paperheight]current page.south west)
      -- ([xshift=0.95\paperwidth, yshift=0.34\paperheight]current page.south west);

    % Contact details
    \node[anchor=north west,
          text=Smoke,
          font=\sffamily\fontsize{8}{14}\selectfont,
          inner sep=0,
          xshift=0.43\paperwidth, yshift=0.31\paperheight]
      at (current page.south west)
      {\shortstack[l]{
        Basado en Melton, A.\\
        Metodología de la Investigación\\
        en Ciencias de la Computación
      }};

  \end{tikzpicture}
\end{frame}
}

\end{document}